\section{Elección de la fuente}

Las fuentes de donde se puede extraer una enzima son muy variadas, pero las principales son:

\begin{itemize}
    \item Vegetales
    \item Animales
    \item Microorganismos
\end{itemize}

\section{Pasos}

\subsection{Obtención de extracto crudo (EC)}
 
En esta etapa se busca liberar la enzima del medio en el que esté, y remover restos celulares insolubles.

Técnicas de bioseparación usadas

\begin{itemize}
    \item Centrifugado
    \item Filtrado: Esta requiere de un cuidado del tamaño del poro
    \item Cromatografía
\end{itemize}

Mediciónes de importancia

\begin{itemize}
    \item Medir la actividad de la enzima
    \item Medir la concentración de la enzima
\end{itemize}

\subsection{Semipurificación - Eliminación de proteínas interferentes}

En esta etapa se busca la semipurificación, esto implica más que nada el eliminar proteínas interferentes que puedan existir en el extracto crudo previamente obtenido.

Los métodos de bioseparación usados en esta etapa suelen ser los siguientes:

\begin{itemize}
    \item Técnicas de bioseparación
    \item Precipitación
    \item Por salado: Este método se apoya del uso de sales para precipitar la proteína. La sal más recomendada es el sulfato de amonio.
    \item Por punto isoelétrico     % REVISAR
\end{itemize}

Dialisis: Este método busca separar las proteínas en base a su peso molecular. Para este método es importante cuidar los siguientes aspectos:

\begin{itemize}
    \item Mantener el extracto a baja temperatura en todo momento.
    \item Realizar cambios de agua constante para eliminar las concentraciones de sales residuales que puedan existir.
\end{itemize}

 Mediciónes de importancia
 
\begin{itemize}
    \item Medir la actividad de la enzima
    \item Medir la concentración de la enzima
\end{itemize}

\subsection{Purificación}

En esta etapa se busca, ahora sí, obtener una proteína pura.
Para esta operación se usan distintos tipos de cromatografía

Técnicas de bioseparación. Cromatografías:

\begin{itemize}
    \item Filtración en gel - Exclusión molecular
    \item Intercambio iónico
    \item Afinidad
\end{itemize}

Medición de importancia: Patrón electroforético